\SourceLeaf{324}%
oder wenn man für \(u\) den Werth aus (3) substituirt, und nach (8), (9), (11),
\[
s=\frac{r\cos\theta^2}{a\cos\varphi}=\frac{r\cos\theta\operatorname{tang}P}{a\operatorname{tang}Q}
\]
setzt
\[
t=\frac{r\cos\theta\operatorname{tang}P\mathord{.}\sin U}{a}(1+\tfrac1{24}uu+\tfrac1{12}ss-\frac{ee\cos P^2}{6\cos\varphi^2}\mathord{.}qq)
\]
Für \(\cos\theta\mathord{.}\operatorname{tang}P=\BareRoot{(1-ee\sin P^2)}\mathord{.}\operatorname{tang}P\) findet man nach (4) den so weit wie hier nöthig ist entwickelten Werth
\[
\BareRoot{(1-ee\sin B^2)}\mathord{.}\operatorname{tang}B(1+\frac{3ee-6e^4\sin P^2+3e^4\sin P^4}{8\cos\varphi^2(1-ee\sin P^2)}\mathord{.}qq)
\]
und folglich
\[
\begin{aligned}
t={}&\frac{rk\operatorname{tang}B\sin U}{a}(1+\frac{5ee+(4ee-14e^4)\sin P^2+5e^4\sin P^4}{24\cos\varphi^2(1-ee\sin P^2)}\mathord{.}qq\\
&\qquad+\tfrac1{12}ss+\tfrac1{24}uu)
\end{aligned}
\]
Macht man nun noch hierin dieselben Substitutionen, wie im vorhergehenden Art., so erhält man als Endresultat (III)
\[
\begin{aligned}
t={}&\frac{rk\operatorname{tang}B\mathord{.}\sin T}{a}(1+\frac{5ee+(4ee-14e^4)\sin B^2+5e^4\sin B^4}{24k^4}\mathord{.}bb\\
&\qquad+\frac{k^4}{12aa\cos\varphi^2}\mathord{.}rr+\tfrac1{24}tt)
\end{aligned}
\]

\indent Die drei Formeln I, II, III enthalten im Wesentlichen die Auflösung unsrer Aufgabe. Dass sie bis zur dritten Ordnung einschliesslich genau sind, steht durch ihre Ableitung unmittelbar fest. Dass aber in der Wirklichkeit ihre Genauigkeit noch eine Ordnung weiter reicht, oder dass der Fehler jeder der Formeln von der fünften Ordnung ist, würde sich leicht durch einige ergänzende Zwischenentwicklungen, oder auch dadurch darthun lassen, dass in den Ausdrücken ihrer Natur nach keine Grössen gerader Ordnung Statt finden können: ich halte mich jedoch dabei nicht auf, da die \textit{zweite} in den folgenden Artikeln (26---32) auszuführende Ableitung der Formeln dasselbe Resultat von selbst in sich begreift.

\Article{26}
Diese Untersuchung ist wie eine selbstständige von allem vorhergehenden unabhängige zu betrachten, und es sollen daher zur Bequemlichkeit und zur Verhütung von Ungewissheiten alle dabei zu verwendenden Bezeichnungen so wie sie auftreten erst erklärt werden. Meistens werden diejenigen Buchstaben, welche schon in der ersten Ableitung gebraucht sind, ihre dortige Bedeutung behalten, doch werden ein Paar derselben (\(u\) und \(s\)), da sie dort bloss Hülfsgrössen vorstel%
